\documentclass[12pt,a4paper]{article}
\usepackage{fontspec}
\usepackage[russian]{babel}
\setmainfont{Liberation Serif}
\usepackage{geometry}
\usepackage{booktabs}
\usepackage{amsmath}
\geometry{left=2.5cm,right=2cm,top=2cm,bottom=2cm}
\setlength{\parskip}{0.4em}
\setlength{\parindent}{1.25em}
\emergencystretch=3em

\begin{document}

\begin{center}
{\Large Лабораторная работа \textnumero~1}\\[0.4em]
{\large Глубокие свёрточные нейронные сети}\\[0.6em]
1 октября 2026~г.
\end{center}

\section{Постановка задачи}

Нужно взять выборку CIFAR-10, собрать свёрточную сеть, обучить её с нуля и посчитать точность классификации на тестовой выборке. В работе рассмотрено пять архитектур. Для максимальной оценки достаточно, чтобы хотя бы одна из них превысила 60\% на тесте; ограничения снизу на точность остальных сетей нет.

\section{Выборка}

CIFAR-10 состоит из 60\,000 цветных изображений размером $32\times32$ пикселя, по 6\,000 на каждый из 10 классов: самолёт, автомобиль, птица, кошка, олень, собака, лягушка, лошадь, корабль, грузовик. Официальное разбиение~--- 50\,000 обучающих изображений и 10\,000 тестовых.

Пиксели переводятся в диапазон $[0,1]$ и нормируются средним и стандартным отклонением по каналам:
\[
\mu = (0{,}4914;\ 0{,}4822;\ 0{,}4465), \qquad
\sigma = (0{,}2470;\ 0{,}2435;\ 0{,}2616).
\]
На обучении к изображению применяются случайное отражение по горизонтали и случайный сдвиг: из кадра с полем ширины 4 пикселя вырезается окно $32\times32$. Тестовые изображения только нормируются.

\section{Алгоритм обучения}

Все пять сетей обучаются одним алгоритмом, с нуля, без предобученных весов.

Функция потерь~--- перекрёстная энтропия между выходом сети и меткой класса. Параметры обновляет Adam: шаг обучения $10^{-3}$, моменты $\beta_1 = 0{,}9$ и $\beta_2 = 0{,}999$, $L_2$-штраф на веса $10^{-4}$. Размер мини-выборки~--- 128 изображений, число эпох~--- 20. Начальные веса свёрток задаются схемой Кайминга, как это делает PyTorch по умолчанию.

Расчёт выполнен в PyTorch 2.9.1+cu128 на графическом процессоре NVIDIA GeForce RTX 3060.

\section{Архитектуры}
\begin{itemize}
\item \textbf{LeNet-подобная.} Две свёртки $5\times5$ (6 и 16 каналов) с ReLU и подвыборкой $2\times2$ после каждой. Полносвязные слои $400 \to 120 \to 84 \to 10$. От классической LeNet-5 отличается функцией ReLU и трёхканальным входом.
\item \textbf{CNN-3.} Три блока <<свёртка $3\times3$ с дополнением, ReLU, подвыборка $2\times2$>> с 32, 64 и 128 каналами. Затем полносвязные слои $2048 \to 256 \to 10$.
\item \textbf{CNN-BN.} Та же схема, что у CNN-3, но после каждой свёртки стоит пакетная нормализация, смещение у свёртки отключено.
\item \textbf{VGG-mini.} Три стадии по две свёртки $3\times3$: 32, 32, подвыборка; 64, 64, подвыборка; 128, 128, подвыборка. Классификатор $2048 \to 256 \to 10$.
\item \textbf{ResCNN.} Ствол $3\to32$, остаточный блок на 32 каналах, свёртка со шагом 2 до 64 каналов, остаточный блок, свёртка со шагом 2 до 128 каналов, остаточный блок, глобальное усреднение и линейный слой $128 \to 10$. В остаточном блоке два свёрточных слоя $3\times3$ с пакетной нормализацией, вход блока прибавляется к выходу.
\end{itemize}

\section{Результаты}

Точность ниже считается на всех 10\,000 тестовых изображениях после последней, двадцатой эпохи. Это точность финальных весов, а не лучшей промежуточной точки.

\begin{center}
\begin{tabular}{lrrr}
\toprule
Архитектура & Параметры & Потеря обучения & Точность теста, \% \\
\midrule
LeNet-подобная & 62\,006 & 1{,}0394 & 66{,}76 \\
CNN-3 & 620\,362 & 0{,}5151 & 80{,}44 \\
CNN-BN & 620\,586 & 0{,}4866 & 81{,}95 \\
VGG-mini & 814\,122 & 0{,}4081 & 84{,}40 \\
ResCNN & 482\,730 & 0{,}2912 & 85{,}78 \\
\bottomrule
\end{tabular}\\[0.4em]
{\small Таблица 1. Результаты пяти свёрточных сетей на CIFAR-10.}
\end{center}

Наибольшая точность у сети <<ResCNN>>: 85{,}78\%. Наименьшая~--- у сети <<LeNet-подобная>>: 66{,}76\%. Порог 60\% превышен у всех пяти сетей.

\section{Вывод}

При одном и том же алгоритме точность определяется архитектурой. По возрастанию: LeNet-подобная~--- 66{,}76\%, CNN-3~--- 80{,}44\%, CNN-BN~--- 81{,}95\%, VGG-mini~--- 84{,}40\%, ResCNN~--- 85{,}78\%. Пакетная нормализация даёт небольшой прирост относительно CNN-3. Две свёртки на стадию и остаточные блоки поднимают результат ещё выше; лучшая сеть~--- ResCNN, и параметров у неё меньше, чем у VGG-mini. Условие на 10 баллов выполнено: архитектур пять, и каждая выше 60\%.

\end{document}
